\begin{center}
    {\large \textbf{Giorno 15} \textbar\ "Kuta is not Ubud" \textbar\ 23 Agosto 2024 \textbar\ \textbf{Kuta}, Bali}
\end{center}

\noindent\rule{\textwidth}{0.4pt}
\begin{figure}[h!] % Portrait images below
    \centering
    % First portrait image (on the left)
    \begin{minipage}[h]{0.48\textwidth} % Set width equal to half the page
        \centering
        \includegraphics[width=\textwidth]{images/flags/Screenshot Kuta Bali.png}
    \end{minipage}
    \hfill
    % Text
    \begin{minipage}[h]{0.48\textwidth}
        \raggedright
        \textbf{Superficie}: 4.567 km² \\
        \vspace{0.3cm}
        \textbf{Popolazione}: 3,3 milioni \\
        \vspace{0.3cm}
        \textit{L'isola}
    \end{minipage}
\end{figure}

\landscapeLargeMargin{images/day15/PXL_20240823_055912832.MP.jpg}

\landscapeLargeMargin{images/day15/PXL_20240823_055921704.MP.jpg}

\landscapeLargeMargin{images/day15/PXL_20240823_035017914.MP.jpg}

\landscapeLargeMargin{images/day15/PXL_20240823_111126598.jpg}

\landscapeLargeMargin{images/day15/PXL_20240823_111558008.MP.jpg}

\landscapeLargeMargin{images/day15/PXL_20240823_111151233.NIGHT.jpg}

\landscapeLargeMargin{images/day15/PXL_20240823_115506224.jpg}

\landscapeLargeMargin{images/day15/PXL_20240823_134937888.MP.jpg}

\noindent 23 Agosto 2024 \newline

La notte nel nuovo hotel procede particolarmente bene, oltre ogni aspettativa. Dormiamo tutta la notte senza essere svegliati e ci alziamo una mezz'ora prima dell'appuntamento con Ketut. Per ottimizzare i tempi dando priorità al sonno, e temendo la qualità della colazione del posto, la sera prima avevamo comprato biscotti e dolciumi da mangiare in macchina durante il tragitto.

Facciamo il check-out semplicemente lasciando le chiavi sul bancone della reception, mentre il ragazzo che ci aveva accolto la sera prima arriva in motorino. Ci dice solo "check-out" e ci salutiamo a gesti.

Questa volta a prelevarci non c'è Ketut, ma un suo collega: il simpatico Nyoman. È originario di Ubud, vive in periferia della città con la moglie e le due figlie di 8 e 15 anni e parla inglese, un po' di spagnolo per lavoro, indonesiano e balinese — che mi spiega essere una lingua completamente a sé e non un dialetto come invece credevo io.

Gli raccontiamo dove siamo stati e gli chiediamo dei mercati del Sukawati. Ci spiega che forse dieci anni fa era ancora possibile trovare arte e artigianato originali dei nativi, ma che negli ultimi anni la città, come anche Ubud, è cambiata molto ed è diventata turistica e commerciale. Sentiamo nella sua voce un lieve senso di dispiacere quando ci racconta com'era Ubud prima del turismo di massa: lui vive fuori dal centro e, oltre al cambiamento della città, racconta che sempre più russi costruiscono ville private vicino al suo villaggio.

Parliamo di cibo italiano — la sua conoscenza si ferma a spaghetti, carbonara e cacio e pepe — di sport, quindi di calcio e MotoGP, di scuola, di lavoro e di famiglia. “Mia moglie fa la maestra, sai”. Mi chiede del movimento LGBT in Italia, per curiosità, ci tiene a precisare, non in maniera discriminatoria. Ricadiamo inevitabilmente anche sul discorso religione e quando gli spiego che in Italia oltre ai cristiani c'è un'alta percentuale di atei sembra immediatamente zittirsi. Non so se ho urtato la sua sensibilità — "esticazzi" in balinese non so come si dica — o se sta solo guardando la strada, ma propendo per la seconda opzione dal momento che sembra uno abbastanza aperto di mentalità e soprattutto perché in quell'istante sbaglia l'uscita e viene intrappolato nel traffico.

Arrivati a destinazione, dopo averci consegnato i bagagli ci saluta e ci dice di fare attenzione perché “Kuta is not Ubud”. Con questa iniezione di fiducia e speranza ci incamminiamo verso l'hotel Truntum, tenendoci ben stretti gli zaini.

Io sono riuscito in queste settimane a fare stare tutto nei miei due zaini, man mano sempre più gonfi. Elisa invece ha iniziato a distribuire le sue cose in sacchetti che tiene appesi a sé come estensione del proprio corpo. Vederla vestita di verde con i sacchetti gialli e rossi che le pendono dalle braccia mi fa pensare a un albero di Natale.

Alla reception del Truntum sono molto gentili: la stanza non è ancora pronta ma mi chiedono il numero di telefono per avvisarmi su WhatsApp quando si fosse liberata. Sto ancora aspettando il messaggio.

\begin{center}
    % Decorative line
    \noindent\rule{0.6\textwidth}{0.4pt}
\end{center}

Lasciamo i bagagli e andiamo a riposarci un po' in piscina. Ormai viaggiando ho imparato a riconoscere il tipo di intrattenimento di una città basandomi sull'aspetto delle persone negli hotel — si tratta di una percezione, ma spesso è azzeccata. La piscina alle undici del mattino è piena di americani in pensione, riconoscibili dal girovita, di gruppetti di giovani e qualche isolata famiglia con bambini che invece di trascorrere la giornata al mare preferisce la piscina dell'albergo. Triste per noi in realtà, perché loro sembrano apprezzare molto il contesto: bar da cui è possibile essere serviti direttamente in acqua, musica house come sottofondo costante e addetti agli asciugamani pronti a soddisfare ogni esigenza.

Se sono tutti qui mi aspetto che le spiagge siano deludenti. Decidiamo comunque di dare una possibilità alla città e andiamo in avanscoperta. In particolare ho messo gli occhi su una libreria di libri usati che spero possa avere la famosa copia in indonesiano di Cent'anni di solitudine.

Camminando per le vie di Kuta il panorama è diverso rispetto a Ubud e alle regioni più a nord. Ci sono sempre i negozi di souvenir con venditori assillanti, ma questa volta sono tutti uomini. Tra i giochi per bambini vediamo per la prima volta pistole e armi giocattolo, probabilmente più adatte ai turisti americani che affollano Kuta e la nostra piscina. Anche locali e ristoranti sembrano allinearsi alla clientela: molte più insegne con pubblicità del Jack Daniel's, bevande in formato super-size — dai frullati agli alcolici — e addirittura un locale che ospita al suo interno un poligono di tiro. Questa cosa la ripeto scandendo bene le parole: un-locale-che-ospita-al-suo-interno-un-poligono-di-tiro. Non ho indagato se le armi fossero ad aria compressa o meno, e non ci tengo particolarmente a saperlo.

Troviamo la libreria e ovviamente anche qui non hanno libri in indonesiano. Perdoname padre, ci ho provato in tutti i modi. Il giro si rivela comunque fruttuoso perché troviamo finalmente le sigarette aromatizzate al caffè e al chiodo di garofano. Inizialmente il mio intento è di comprarne più pacchetti come souvenir, poi ne assaggio una — sì, buona, ma non fumo — e siamo contenti così.

Ci avventuriamo poi sulla strada principale che costeggia la spiaggia: visitiamo un centro commerciale fronte mare che è forse la cosa più occidentale vista in queste due settimane e usciamo subito. Eli vorrebbe andare a mangiare e poi tornare nel pomeriggio in spiaggia, ma insisto per passarci prima, così da toglierci subito il dubbio invece di tornare poi a vuoto. Si rivela la scelta giusta.

La spiaggia è terribile. Il mare è molto agitato con onde alte adatte solo ai surfisti — probabilmente tutti provenienti dalla vicina Australia — che infatti sono gli unici in acqua. Il bagnasciuga si divide su tre livelli: una prima fila più vicina al mare con lettini e ombrelloni, una seconda fascia con ombrelloni e sedie di plastica su cui sedersi e probabilmente bere birra — ho già detto che ci sono turisti americani qui? — e una terza zona piena di bancarelle e bar improvvisati tra cui camminare per essere importunati da chi vuole venderti souvenir, alcol, un posto sul lettino o droga. Ci viene in mente la frase di Nyoman sul fare attenzione, e realizziamo che i volti sorridenti degli indonesiani conosciuti a Ubud e a Lombok sono qui sostituiti da facce da criminali. Non ho mai avuto nulla contro i tatuaggi ma qui per la prima volta li vediamo sui locali, lo stesso dicasi per i capelli lunghi sporchi e spettinati.

La passeggiata è sufficiente per decidere di trascorrere il resto del pomeriggio in piscina. Prima però dobbiamo pensare al pranzo e, essendo stanchi, non abbiamo voglia di rischiare. Con tutto questo turismo americano ci deve per forza essere un McDonald's e in effetti è esattamente dietro al nostro albergo. La meta per il pranzo è scelta: McChicken sotto un getto d'aria condizionata perfetto per mantenere a temperatura l'azoto liquido. Oziamo un po' nel locale cercando di buttar giù le patatine, che qui servono solo con salsa locale piccantissima, poi torniamo in hotel e andiamo direttamente in piscina.

Come al mattino c'è il pienone e inizialmente trova posto solo Elisa. Mentre la accompagno passo a fianco a due signore americane obese che mangiano hamburger e patatine sui loro lettini. Faccio quattro chiacchiere col ragazzo degli asciugamani per inquadrare un po' la situazione, e tutto ciò che mi dice va contro quello che ho imparato in queste settimane: il suo focus sono i party e mi chiede se sono stato nelle città festaiole di Bali. Quando capisce che non ne abbiamo visitata neanche una e che soprattutto non ho un particolare interesse nel farlo, mi consiglia di andare a Nusa Penida, che a suo dire è splendida.

Mi viene in mente che stamattina con Nyoman avevamo parlato proprio di Nusa Penida: il suo commento era stato che anche quest'isola ormai è diventata molto commerciale, che in questa stagione le mante tendono a rimanere in fondo al mare e che l'unico modo per vederle è acquistare tour in barca con capitani particolarmente esperti. Enrico e Cristina ci sono stati e ci hanno detto di aver fatto un semplice giro di snorkeling senza vedere alcuna manta. Quando racconto al ragazzo che a Gili Trawangan abbiamo visto le tartarughe, si affretta a dire che è possibile vederle anche qui — però precisa che si tratta di tartarughine appena nate, liberate in mare di tanto in tanto come attrazione turistica. Paragonato al nuotare accanto a tartarughe giganti, è un po' come fare snorkeling in una vasca da bagno.

Quando si libera un posto raggiungo Elisa e cerco di rimanere all'ombra dell'ombrellone finché riesco; non appena mi ritrovo tutto al sole decidiamo di andare in acqua. Alla reception ci avevano dato un buono per una consumazione gratuita, così lo presentiamo al bancone del bar dove ci comunicano che possiamo usarlo liberamente solo al ristorante dell'albergo, e che usandolo in piscina l'unica cosa disponibile sarebbe il succo d'arancia. Questo cercare di vincolarci nelle scelte ha lo stesso effetto dei venditori che cercano di importi i souvenir — ci allontana soltanto. Va sereno e dacci ste due aranciate, capo!

\begin{center}
    % Decorative line
    \noindent\rule{0.6\textwidth}{0.4pt}
\end{center}

Quando finalmente arriva l'ora del check-in vado a ritirare le chiavi e salgo in stanza. Si tratta di una camera d'albergo standard, molto grande e spaziosa con i classici comfort — perfetta per un viaggio aziendale, insomma. Mi corico un attimo a riposare: la musica costante della piscina mi ha fatto venire mal di testa.

Appena Eli si sveglia mi raggiunge e iniziamo a preparare le valigie per l'indomani. Ho seriamente paura che Elisa non riesca a far stare tutto: di souvenir decenti non ne abbiamo trovati molti, ma senza che me ne accorgessi durante la vacanza — laddove ha potuto — ha portato via conchiglie, addobbi e tutto ciò che negli hotel non era incatenato. In totale: cinque kit monouso di spazzolino e dentifricio, quattro mascherine per dormire, otto cuffie da doccia, un numero indefinito di bustine di tè rubate in hotel con altrettante bustine di zucchero — che giustamente a casa ci manca — conchiglie a sufficienza da riempire ogni mensola di casa e tanti, tanti, ma tanti cotton fioc.

Incredibilmente riesce a far entrare tutto, fatta eccezione per un ananas, che tocca mettere nel mio zaino. Farò tutto il viaggio di ritorno con un ananas nello zaino. Terminata l'operazione facciamo varie congetture su cosa potrebbero dirci in aeroporto riguardo alla scatola di cartone coi souvenir di terracotta di Lombok, ma decidiamo di presentarci senza modificare il bagaglio.

\begin{center}
    % Decorative line
    \noindent\rule{0.6\textwidth}{0.4pt}
\end{center}

Quando usciamo scopriamo che sul retro dell'hotel c'è un'altra piscina e che è possibile cenarvi attorno, ma per l'ultima cena vogliamo qualcosa di più tipico. Ripercorriamo le strade del mattino evitando il più possibile i venditori — non è facile — e finalmente troviamo un locale all'aperto molto carino che serve piatti tradizionali. Via di nuovo col menù collaudato: Mie Goreng, spring rolls, Bintang e Bintang Radler e per dolce dei pancake, che si rivelano poi essere crepe alla banana.

Ci godiamo l'atmosfera del locale: niente musica, luci soffuse, in mezzo al verde. Siamo seduti sotto una pianta di manghi e verifichiamo bene di non essere a tiro di frutti particolarmente maturi.

Usciti ci avviamo verso la spiaggia e durante il tragitto veniamo continuamente interpellati da locali che chiedono se abbiamo bisogno di droga. Li ignoriamo, ma mi viene in mente che sarebbe divertente prenderli in contropiede e rispondergli con la faccia più seria possibile: "No drugs, but can you find me some giant turtles?" Questa scena non si avvera, purtroppo.

Approfittiamo della passeggiata per le vie interne per comprare qualche braccialetto di pietra lavica e minerali — giada, agata, quarzo azzurro — stavolta contrattando con tenacia. Un plauso particolare a Didì, un venditore che alla richiesta di un braccialetto più largo si è messo d'impegno con ago e filo e me l'ha realizzato su misura sul momento.

Terminati gli acquisti ci dirigiamo sulla strada che costeggia il mare, passando davanti ai locali dove assistiamo a uno spettacolo di un barman che fa il giocoliere con bottiglie e bicchieri da cocktail e a musicisti live decisamente più stonati che a Ubud. Ci fermiamo in un localino tranquillo dove ci godiamo una birra e ripensiamo ai momenti migliori di queste settimane. Il primo pensiero va ovviamente a Oshi e al soggiorno nel Tetebatu a Lombok, sicuramente l'esperienza migliore vissuta da entrambi. Se lo dico nonostante il Muezzin alle quattro e la privazione del sonno costante, sono decisamente onesto.

Finite le birre torniamo in hotel dove tutto è già pronto per l'indomani e crolliamo addormentati.

\begin{center}
    % Decorative line
    \noindent\rule{0.6\textwidth}{0.4pt}
\end{center}

\textit{P.s. Oltre all'ananas scoprirò solo in un secondo momento di aver viaggiato con quattro salak — i frutti simili al fico che neanche le piacciono — ma che deve portarsi a casa a tutti i costi.}


\pagestyle{empty} % disable numbers

\landscapeThinMargin{images/day15/PXL_20240823_135010971.PORTRAIT.jpg}

\landscapeThinMargin{images/day15/PXL_20240823_031257392.MP.jpg}

\pagestyle{fancy} % enable numbers